\section{Operador constitutivo y reconstrucción de los canales del vacío}
\label{iii:sec:constitutiva}

Las magnitudes constitutivas se obtienen mediante una operación sobre el
transporte angular. La información de partida comprende los canales, sus
proyectores y la orientación que los distingue. Esta organización permite
seguir por separado la construcción del operador, la evaluación de sus
funciones espectrales y su realización dimensional posterior.

\subsection{Respuesta positiva sobre dos hojas}

En la cámara orientada $x>y>0$ de la sección~\ref{iii:sec:hojas}, sean
$P_\pm=(I\pm\mathcal R)/2$ y
$T=q_+P_++q_-P_-$, con $0<q_+<q_-<1$. Las incidencias construidas por el
calendario fijan los sectores $90$ y $120$. La respuesta constitutiva es
\begin{equation}
 \mathcal V=(I-T^{90})(I-T^{120})^{-1}.
 \label{iii:eq:vacuum-op}
\end{equation}
Esta definición recibe el transporte procedente del TPK y conserva las
asignaciones de hoja. La regla de respuesta y su dominio son parte explícita
de la construcción.

El cociente procede de comparar las dos incidencias del calendario sobre
un mismo transporte, no de ajustar por separado cuatro magnitudes. Si
$S=T^{30}$ y $\Sigma_j(S)=I+S+\cdots+S^{j-1}$, la factorización de
$I-S^j$ da
\begin{equation}
 \mathcal V\,\Sigma_4(S)=\Sigma_3(S),\qquad
 \mathcal V=\Sigma_3(S)\Sigma_4(S)^{-1}.
 \label{iii:eq:completion-response}
\end{equation}
Los factores $\Sigma_3$ y $\Sigma_4$ expresan respectivamente los
sectores $90=3\cdot30$ y $120=4\cdot30$. Como el espectro de $S$ está
contenido en $(0,1)$, $\Sigma_4(S)$ es invertible: la ecuación de
completación determina una única respuesta para el transporte fijado.
Esta unicidad concierne a la regla constitutiva declarada y a sus
incidencias. La identificación electromagnética añade la asignación
orientada de las hojas y las rectas dimensionales que se precisan más
abajo.

\begin{proposition}[Descomposición y positividad]
La respuesta está bien definida y
\begin{equation}
 \mathcal V=r_+P_++r_-P_-,\qquad
 r_\pm=\frac{1-q_\pm^{90}}{1-q_\pm^{120}}.
 \label{iii:eq:rchannels}
\end{equation}
Sus autovalores satisfacen $3/4<r_-<r_+<1$ en esta cámara.
\end{proposition}
\begin{proof}
Cada coeficiente $1-q_\pm^{120}$ es positivo. La inversa de $I-T^{120}$
es la suma de los proyectores divididos por esos coeficientes, lo que prueba
\eqref{iii:eq:rchannels}. Con $s=q^{30}$ se obtiene
\begin{equation}
 f(s)=\frac{1+s+s^2}{1+s+s^2+s^3},\qquad
 f'(s)=-\frac{s^2(3+2s+s^2)}{(1+s+s^2+s^3)^2}<0.
 \label{iii:eq:monotone}
\end{equation}
Los límites en $0$ y $1$ son $1$ y $3/4$, respectivamente. La monotonía
y el orden de los canales concluyen la prueba.
\end{proof}

\subsection{Cuatro coordenadas de una respuesta común}

Se definen las evaluaciones normalizadas
\begin{equation}
 \widehat\varepsilon=r_+^2,\qquad
 \widehat\mu=r_-^2,\qquad
 \widehat Z=\frac{r_-}{r_+},\qquad
 \widehat c=\frac1{r_+r_-}.
 \label{iii:eq:four}
\end{equation}
La primera pareja conserva las respuestas cuadráticas asignadas a cada hoja.
El producto recoge el modo común; el cociente distingue su orientación relativa.
Se verifican exactamente
\begin{equation}
 \widehat\mu\widehat\varepsilon=\widehat c^{-2},\qquad
 \frac{\widehat\mu}{\widehat\varepsilon}=\widehat Z^2,
 \qquad
 r_+=\frac1{\sqrt{\widehat Z\widehat c}},\quad
 r_-=\sqrt{\frac{\widehat Z}{\widehat c}}.
 \label{iii:eq:four-identities}
\end{equation}
Las raíces positivas fijan la inversión. Bajo la conjugación de hojas,
leída respecto de las marcas $P_\pm$ fijas, se
intercambian $\widehat\varepsilon$ y $\widehat\mu$, se invierte
$\widehat Z$ y se conserva $\widehat c$. La respuesta se transporta entonces
a la cámara $x>-y>0$: las fórmulas siguen siendo válidas, con los órdenes
$q_-<q_+$ y $r_+<r_-$ invertidos. Por consiguiente, el modo común
y la orientación son datos recuperables mediante observables complementarios.

\begin{proposition}[Determinación de las evaluaciones cuadráticas]
Sea $W=r_+P_++r_-P_-$ una respuesta positiva sobre las dos hojas.
Fijadas las lecturas de producto y cociente
\[
 \widehat c^{-1}=\det W=r_+r_-,
 \qquad \widehat Z=r_-/r_+,
\]
existe una única pareja positiva $(\widehat\varepsilon,\widehat\mu)$
que satisface
$\widehat\mu\widehat\varepsilon=\widehat c^{-2}$ y
$\widehat\mu/\widehat\varepsilon=\widehat Z^2$.
Es precisamente $(r_+^2,r_-^2)$.
\end{proposition}
\begin{proof}
La positividad permite tomar las raíces sin ambigüedad:
\[
 \widehat\mu=\frac{\widehat Z}{\widehat c}
 =\frac{r_-}{r_+}r_+r_-=r_-^2,\qquad
 \widehat\varepsilon=\frac1{\widehat Z\widehat c}
 =\frac{r_+}{r_-}r_+r_-=r_+^2.
\]
Las expresiones verifican las dos relaciones y son las únicas raíces
positivas de ellas.
\end{proof}

Así, los cuadrados no constituyen dos elecciones independientes añadidas
a los canales: quedan determinados por las lecturas multiplicativa y
orientada, una vez fijadas las relaciones constitutivas. Esta proposición
establece la unicidad dentro de ese sistema de reglas. La interpretación
de las reglas como respuesta electromagnética conserva su contenido
adicional de representación física.

La lectura determinantal del transporte y la respuesta constitutiva
también deben conservar su orden. Para $T=e^{-xI-y\mathcal R}$ se tiene
$\det T=e^{-2x}$, mientras que
\[
 \det\mathcal V=f(q_+^{30})f(q_-^{30})=\widehat c^{-1}.
\]
Ambas proceden del mismo operador de dos hojas, pero evalúan funciones
diferentes de su espectro. En general, aplicar primero el determinante
y después una función racional no equivale a tomar el determinante de
esa función del operador. Mantener los dos autovalores hasta completar
la respuesta conserva la información de orientación que un determinante
aislado ya no distingue.

\subsection{Inversión cúbica y recuperación angular}

\begin{theorem}[Inversión constitutiva en el dominio contractivo]
Para cada $r\in(3/4,1)$ existe una única $s\in(0,1)$ tal que $f(s)=r$.
Esta $s$ es la raíz única en dicho intervalo de
\begin{equation}
 r s^3+(r-1)(1+s+s^2)=0.
 \label{iii:eq:cubic}
\end{equation}
La respuesta etiquetada reconstruye $s_\pm$, $q_\pm=s_\pm^{1/30}$ y
\begin{equation}
 x=-\frac12\log(q_+q_-),\qquad
 y=\frac12\log\frac{q_-}{q_+}.
 \label{iii:eq:recover-angular}
\end{equation}
\end{theorem}
\begin{proof}
La derivada y los límites de \eqref{iii:eq:monotone} prueban que $f$ es una
biyección decreciente entre los intervalos indicados. Multiplicar $f(s)=r$
por su denominador produce \eqref{iii:eq:cubic}. La raíz positiva de orden
$30$ recupera el canal. Finalmente, suma y diferencia de
$\log q_\pm=-x\mp y$ dan \eqref{iii:eq:recover-angular}.
\end{proof}

La reconstrucción conserva las etiquetas de hoja. Con el operador completo
$\mathcal V$ se conservan también los proyectores y se aplica la función
inversa por cálculo espectral. Recuperar dos coordenadas angulares es una
operación definida sobre esta representación; la traza histórica completa
del TPK conserva datos adicionales.

El intervalo $r\in(3/4,1)$ y la función inversa de este teorema pertenecen
a la normalización interna de \eqref{iii:eq:four}. Si se dispone de
coordenadas expresadas en otra carta de unidades, se recupera primero
esta normalización mediante la transformación de la
sección~\ref{iii:sec:evaluacion}; después se aplica la inversión cúbica.

\subsection{Coordenadas logarítmicas y refinamiento}

Sean $\mathcal E=\log(r_+r_-)$ y $\mathcal B=\log(r_-/r_+)$. Entonces
\begin{equation}
 \log\mathcal V=\tfrac12\mathcal E I-\tfrac12\mathcal B\mathcal R.
 \label{iii:eq:logvac}
\end{equation}
Esta descomposición exhibe la componente invariante y la componente orientada.
En la amplificación $\mathcal V_m=\mathcal V\otimes I_{9^m}$, las
expectativas parciales normalizadas satisfacen
$E_m(\mathcal V_{m+1})=\mathcal V_m$. Para todo polinomio $p$,
$p(\mathcal V_m)=p(\mathcal V)\otimes I_{9^m}$; de ahí
$E_m(p(\mathcal V_{m+1}))=p(\mathcal V_m)$. La continuidad uniforme
extiende la identidad a las funciones continuas del espectro. Se conservan
así los proyectores, las potencias y el logaritmo en la torre indicada.

\subsection{Tipos dimensionales de las cuatro evaluaciones}

Con las rectas de acción, carga y velocidad, provistas de secciones
$u_S,u_Q,u_C$, la realización dimensional se escribe
\begin{align}
 Z_{\rm vac}&=\widehat Z u_Su_Q^{-2},&
 c_{\rm vac}&=\widehat c u_C,\\
 \mu_{\rm vac}&=\widehat\mu u_Su_C^{-1}u_Q^{-2},&
 \varepsilon_{\rm vac}&=\widehat\varepsilon u_S^{-1}u_C^{-1}u_Q^2.
 \label{iii:eq:dimensions}
\end{align}
Multiplicando y dividiendo estas expresiones se preservan
$\mu_{\rm vac}\varepsilon_{\rm vac}=c_{\rm vac}^{-2}$ y
$\mu_{\rm vac}/\varepsilon_{\rm vac}=Z_{\rm vac}^2$.
La carta SI evalúa después las secciones. La comparación física requiere
exponer su construcción y normalización junto a la evaluación; las cuatro
coordenadas normalizadas de \eqref{iii:eq:four} conservan su tipo adimensional.
